\documentclass[12pt,a4paper]{article}
\usepackage[utf8]{inputenc}
\usepackage[ngerman]{babel}
\usepackage{amsmath}
\usepackage{booktabs}
\usepackage{geometry}
\usepackage{listings}
\usepackage{xcolor}
\usepackage{hyperref}

\geometry{margin=2.5cm}

\title{Gehaltsberechnung für Kandidaten\\
\large Dokumentation der \texttt{fill\_salary.py} Logik}
\author{KI-unterstütztes BMS}
\date{\today}

\begin{document}

\maketitle
\tableofcontents
\newpage

\section{Übersicht}

Das Skript \texttt{fill\_salary.py} berechnet den Gehaltswunsch für Kandidaten in der PostgreSQL-Datenbank basierend auf drei Hauptfaktoren:

\begin{itemize}
    \item \textbf{Position} (z.B. Assistenzarzt, Facharzt, Oberarzt, Chefarzt)
    \item \textbf{Fachbereich} (z.B. Radiologie, Innere Medizin, Chirurgie)
    \item \textbf{Land} (z.B. Deutschland, Schweiz, Österreich)
\end{itemize}

\section{Berechnungsformel}

Das Gehalt wird nach folgender Formel berechnet:

\begin{equation}
    G_{\text{final}} = \text{round}\left(\left(G_{\text{basis}} \times F_{\text{Fach}} \times F_{\text{Land}} + V + Z\right) / 500\right) \times 500
\end{equation}

wobei:
\begin{align*}
    G_{\text{basis}} &= \text{Basisgehalt der Position} \\
    F_{\text{Fach}} &= \text{Fachbereich-Faktor} \\
    F_{\text{Land}} &= \text{Land-Faktor} \\
    V &= \text{Varianz} \in [-0.15 \times G, +0.15 \times G] \\
    Z &= \text{Zufallsstreuung} \in [-2000, +2000] \text{ EUR}
\end{align*}

Die Rundung auf 500 EUR sorgt für realistische Gehaltsstufen.

\section{Positionshierarchie}

\begin{table}[h]
\centering
\begin{tabular}{@{}llr@{}}
\toprule
\textbf{Position} & \textbf{Level} & \textbf{Basisgehalt (EUR)} \\
\midrule
Assistenzarzt & 1 & 60\,000 \\
Facharzt & 2 & 100\,000 \\
Oberarzt & 3 & 140\,000 \\
Leitender Oberarzt & 4 & 170\,000 \\
Chefarzt & 5 & 230\,000 \\
Standortleiter & 6 & 200\,000 \\
Gesellschafter & 7 & 280\,000 \\
\bottomrule
\end{tabular}
\caption{Basisgehälter pro Position in Deutschland}
\label{tab:positionen}
\end{table}

\section{Fachbereich-Faktoren}

Die Fachbereich-Faktoren spiegeln die unterschiedlichen Gehaltsniveaus verschiedener medizinischer Spezialisierungen wider:

\begin{table}[h]
\centering
\begin{tabular}{@{}lr@{}}
\toprule
\textbf{Fachbereich} & \textbf{Faktor} \\
\midrule
Radiologie & 1.70 \\
Nuklearmedizin & 1.65 \\
Neuroradiologie & 1.68 \\
Strahlentherapie & 1.60 \\
Kinderradiologie & 1.55 \\
Mammographie & 1.50 \\
Anästhesie & 1.30 \\
Chirurgie & 1.25 \\
Orthopädie \& UCH & 1.25 \\
Innere Medizin & 1.15 \\
Gynäkologie & 1.15 \\
Pädiatrie/Kindermedizin & 1.10 \\
Psychiatrie & 1.00 \\
\bottomrule
\end{tabular}
\caption{Fachbereich-Faktoren (relativ)}
\label{tab:fachbereiche}
\end{table}

\textbf{Hinweis:} Radiologie und bildgebende Verfahren haben die höchsten Faktoren, da diese Fachbereiche im Markt höhere Gehälter zahlen.

\section{Länder-Faktoren}

Die Länder-Faktoren berücksichtigen internationale Gehaltsunterschiede:

\begin{table}[h]
\centering
\begin{tabular}{@{}lrr@{}}
\toprule
\textbf{Land} & \textbf{Faktor} & \textbf{Relativ zu DE (\%)} \\
\midrule
Schweiz & 1.50 & +50\% \\
Luxemburg & 1.35 & +35\% \\
Dänemark & 1.20 & +20\% \\
Niederlande & 1.05 & +5\% \\
Belgien & 1.03 & +3\% \\
Deutschland & 1.00 & Basis \\
Frankreich & 0.98 & -2\% \\
Österreich & 0.95 & -5\% \\
Tschechien & 0.45 & -55\% \\
Polen & 0.40 & -60\% \\
\bottomrule
\end{tabular}
\caption{Länder-Faktoren (relativ zu Deutschland)}
\label{tab:laender}
\end{table}

\section{Datenextraktion}

\subsection{Position-Erkennung}

Die Position wird aus zwei Datenbank-Feldern extrahiert:
\begin{enumerate}
    \item \textbf{Primär:} \texttt{position\_now} -- aktuelle Position
    \item \textbf{Fallback:} \texttt{qualification} -- Qualifikation
    \item \textbf{Standard:} Falls keine Position erkannt wird $\rightarrow$ ``Facharzt''
\end{enumerate}

\subsection{Fachbereich-Erkennung}

Der Fachbereich wird aus dem Feld \texttt{department} extrahiert:
\begin{itemize}
    \item Text-Normalisierung (Kleinschreibung, Trimming)
    \item Pattern-Matching gegen Fachbereich-Liste
    \item \textbf{Fallback:} ``Innere Medizin'' (mittlerer Gehaltsfaktor)
\end{itemize}

\subsection{Land-Erkennung}

Das Land wird hierarchisch aus folgenden Feldern extrahiert:
\begin{enumerate}
    \item \textbf{Primär:} \texttt{wunscharbeitsort}
    \item \textbf{Sekundär:} \texttt{wohnort}
    \item \textbf{Tertiär:} \texttt{arbeitsort}
    \item \textbf{Standard:} Deutschland
\end{enumerate}

\textbf{Länder-Mapping:} Verschiedene Schreibweisen werden normalisiert:
\begin{itemize}
    \item Deutschland: ``deutschland'', ``de'', ``ger'', ``germany''
    \item Schweiz: ``schweiz'', ``ch'', ``switzerland''
    \item Österreich: ``österreich'', ``oesterreich'', ``at'', ``austria''
\end{itemize}

\section{Rechenbeispiele}

\subsection{Beispiel 1: Oberarzt Radiologie in Deutschland}

\begin{align*}
    G_{\text{basis}} &= 140\,000 \text{ EUR (Oberarzt)} \\
    F_{\text{Fach}} &= 1.70 \text{ (Radiologie)} \\
    F_{\text{Land}} &= 1.00 \text{ (Deutschland)} \\
    G &= 140\,000 \times 1.70 \times 1.00 = 238\,000 \text{ EUR} \\
    \text{Mit Varianz:} &\quad V \in [-35\,700, +35\,700] \\
    \text{Mit Streuung:} &\quad Z \in [-2\,000, +2\,000] \\
    G_{\text{final}} &\approx 200\,000 \text{ bis } 275\,000 \text{ EUR}
\end{align*}

\subsection{Beispiel 2: Facharzt Radiologie in der Schweiz}

\begin{align*}
    G_{\text{basis}} &= 100\,000 \text{ EUR (Facharzt)} \\
    F_{\text{Fach}} &= 1.70 \text{ (Radiologie)} \\
    F_{\text{Land}} &= 1.50 \text{ (Schweiz)} \\
    G &= 100\,000 \times 1.70 \times 1.50 = 255\,000 \text{ EUR} \\
    G_{\text{final}} &\approx 215\,000 \text{ bis } 295\,000 \text{ EUR}
\end{align*}

\subsection{Beispiel 3: Assistenzarzt Psychiatrie in Deutschland}

\begin{align*}
    G_{\text{basis}} &= 60\,000 \text{ EUR (Assistenzarzt)} \\
    F_{\text{Fach}} &= 1.00 \text{ (Psychiatrie)} \\
    F_{\text{Land}} &= 1.00 \text{ (Deutschland)} \\
    G &= 60\,000 \times 1.00 \times 1.00 = 60\,000 \text{ EUR} \\
    G_{\text{final}} &\approx 49\,000 \text{ bis } 71\,000 \text{ EUR}
\end{align*}

\section{Algorithmus-Ablauf}

\begin{enumerate}
    \item \textbf{Datenbankverbindung herstellen}
    \item \textbf{Spalte prüfen/erstellen:} Falls \texttt{gehaltswunsch} nicht existiert, wird sie als \texttt{INTEGER} hinzugefügt
    \item \textbf{Kandidaten laden:} Alle Kandidaten mit relevanten Feldern werden geladen
    \item \textbf{Für jeden Kandidaten:}
    \begin{enumerate}
        \item Position aus \texttt{position\_now} oder \texttt{qualification} extrahieren
        \item Fachbereich aus \texttt{department} extrahieren
        \item Land aus \texttt{wunscharbeitsort}/\texttt{wohnort}/\texttt{arbeitsort} extrahieren
        \item Gehalt nach Formel berechnen
        \item \texttt{UPDATE candidates SET gehaltswunsch = ...}
    \end{enumerate}
    \item \textbf{Commit und Statistiken ausgeben}
\end{enumerate}

\section{Statistiken und Validierung}

Das Skript erstellt folgende Statistiken:
\begin{itemize}
    \item Anzahl aktualisierter Kandidaten
    \item Anzahl überschriebener Werte
    \item Kandidaten mit/ohne erkannte Position
    \item Kandidaten mit/ohne erkannten Fachbereich
    \item Länder-Verteilung
    \item Gehaltsverteilung nach Position (Min, Durchschnitt, Max)
\end{itemize}

\textbf{Beispiel-Ausgabe:}
\begin{verbatim}
Oberarzt: 156 Kandidaten | Min: 95,000€ | Ø: 152,500€ | Max: 215,000€
Facharzt: 842 Kandidaten | Min: 72,000€ | Ø: 118,000€ | Max: 175,000€
\end{verbatim}

\section{Besonderheiten}

\subsection{Varianz und Streuung}

Die Kombination aus prozentualer Varianz (±15\%) und absoluter Streuung (±2000 EUR) erzeugt realistische Gehaltsverteilungen:
\begin{itemize}
    \item \textbf{Prozentuale Varianz:} Höhere Positionen haben größere absolute Schwankungen
    \item \textbf{Absolute Streuung:} Fügt feingranulare Unterschiede hinzu
    \item \textbf{Rundung auf 500 EUR:} Erzeugt typische Gehaltsstufen
\end{itemize}

\subsection{Überschreiben bestehender Werte}

Das Skript überschreibt \textbf{alle} Gehaltswerte, auch bereits vorhandene. Dies ist gewünscht, um:
\begin{itemize}
    \item Fehlerhafte alte Berechnungen zu korrigieren
    \item Aktualisierte Faktoren anzuwenden
    \item Konsistenz in der gesamten Datenbank zu gewährleisten
\end{itemize}

\section{Verwendung}

\subsection{Ausführung}

\begin{verbatim}
python src/PostreSQL/import/fill_salary.py
\end{verbatim}

\subsection{Beispiel-Modus}

Das Skript zeigt zunächst Beispielberechnungen:
\begin{verbatim}
=== Beispiel-Gehaltsberechnungen ===

Oberarzt in Radiologie (Deutschland): 238,500€
Facharzt in Radiologie (Deutschland): 170,000€
Oberarzt in Innere medizin (Deutschland): 161,000€
...
\end{verbatim}

Danach wird eine Bestätigung abgefragt, bevor die Datenbank aktualisiert wird.

\section{Datenbank-Schema}

Die relevanten Spalten in der \texttt{candidates}-Tabelle:

\begin{table}[h]
\centering
\begin{tabular}{@{}lll@{}}
\toprule
\textbf{Spalte} & \textbf{Typ} & \textbf{Verwendung} \\
\midrule
\texttt{id} & INTEGER & Primärschlüssel \\
\texttt{position\_now} & VARCHAR & Position-Erkennung (primär) \\
\texttt{qualification} & VARCHAR & Position-Erkennung (fallback) \\
\texttt{department} & VARCHAR & Fachbereich-Erkennung \\
\texttt{wunscharbeitsort} & VARCHAR & Land-Erkennung (primär) \\
\texttt{wohnort} & VARCHAR & Land-Erkennung (sekundär) \\
\texttt{arbeitsort} & VARCHAR & Land-Erkennung (tertiär) \\
\texttt{gehaltswunsch} & INTEGER & Berechnetes Gehalt (Ziel) \\
\bottomrule
\end{tabular}
\caption{Relevante Datenbank-Spalten}
\label{tab:db_schema}
\end{table}

\section{Zusammenfassung}

Die Gehaltsberechnung basiert auf einem dreistufigen multiplikativen Modell mit zusätzlicher Varianz für realistische Streuung. Die Hierarchie der Positionen, die Marktbedingungen der Fachbereiche und internationale Gehaltsunterschiede werden systematisch berücksichtigt.

Das Ergebnis sind plausible, differenzierte Gehaltswünsche für alle Kandidaten im System, die als Grundlage für Matching-Algorithmen und Gehaltsverhandlungen dienen.

\end{document}
